% ! TEX root = ../am2-18-19.tex

\chapter{Stabilitate și linii de cîmp}

\indent\indent Problema stabilității soluțiilor este una foarte importantă și dificilă,
în general. Ideea de bază este că, dacă ne gîndim la interpretarea fizică
a sistemelor de ecuații și soluțiilor lor, de exemplu în cazul mecanicii,
soluția returnată este posibil să fie validă într-o vecinătate foarte mică
a punctului în care a fost calculată. În cazul acesta, situația corespunde
unei poziții de \textit{echilibru (mecanic) instabil}, adică atunci cînd
deplasarea corpului la o mică distanță de poziția curentă îl face să iasă
din poziția de echilibru. Similar, putem avea și \textit{echilibru (mecanic) stabil}
sau \textit{echilibru (mecanic) indiferent}.

Teoria stabilității soluțiilor sistemelor diferențiale este un subiect amplu,
fondat și discutat în detaliu de \textit{H.\ Poincar\'e} și \textit{A.\ Liapunov},
de aceea, majoritatea teoremelor și rezultatelor importante le sînt atribuite.

Contextul în care lucrăm este următorul. fie un sistem diferențial de forma
$ x^\prime = v(t, x) $, unde $ x $ poate fi vector (adică să țină loc de mai
multe variabile). Presupunem că sînt îndeplinite condițiile teoremei fundamentale
de existență și unicitate a soluției problemei Cauchy, pentru $ t \in [t_0, \infty) $
și că avem $ x \in U \seq \RR^n $, un deschis. Așadar, pentru orice $ x_0 \in U $,
există și este unică o soluție $ x = \varphi(t) $, cu $ \varphi : [t_0, \infty) \to U $<
astfel încît $ \varphi(t_0) = x_0 $.

Tipurile de stabilitate sînt definite mai jos.

\begin{definition}\label{def:stabil-inf}
  \index{stabilitate!spre $\infty$}
  O soluție $ x = \varphi(t) $ ca mai sus se numește \emph{stabilă spre $\infty$}
  în sens Poincar\'e-Lia\-pu\-nov (sau, echivalent, $ x_0 = \varphi(t_0) $
  se numește \emph{poziție de echilibru}) dacă la variații mici ale
  lui $ x_0 $ ob\-ți\-nem variații mici ale soluției.

  Formal, pentru orice $ \varepsilon > 0 $, există $ \delta(\varepsilon) > 0 $
  astfel încît pentru orice $ \tilde{x_0} \in \RR^n $, cu
  $ ||\tilde{x_0} - x_0 || < \delta(\varepsilon) $, soluțiile corespunzătoare,
  $ \tilde{\varphi}(t) $ și $ \varphi(t) $ satisfac inegalitatea:
  \[
    | \tilde{\varphi}(t) - \varphi(t) | < \varepsilon, \quad \forall t \geq t_0.
  \]
\end{definition}

\begin{definition}\label{def:stabil}
  \index{stabilitate!Poincar\'e-Liapunov}
  Poziția de echilibru $ x = 0 $ se numește \emph{stabilă în sens Poincar\'e-Liapunov}
  dacă pentru orice $ \varepsilon > 0 $, există $ \delta > 0 $, care depinde
  doar de $ \varepsilon $, astfel încît pentru orice $ x_0 \in U $ pentru care
  $ ||x_0|| < \delta $, soluția $ \varphi $ a sistemului cu condiția inițială
  $ \varphi(0) = x_0 $ se prelungește pe întreaga semiaxă $ t > 0 $ și
  satisface inegalitatea $ || \varphi(t) || < \varepsilon, \forall t > 0 $.

\end{definition}

\begin{definition}\label{def:stabil-as}
  \index{stabilitate!asimptotică}
  Poziția de echilibru $ x = 0 $ a sistemului diferențial autonom se
  numește \emph{asimptotic stabilă} dacă este stabilă și, în plus, pentru
  soluția $ \varphi(t) $ din definiția de mai sus, are loc:
  \[
    \lim_{t \to \infty} \varphi(t) = 0.
  \]
\end{definition}

În exercițiile pe care le vom întîlni, vom folosi următorul rezultat fundamental.

Presupunem că avem un sistem diferențial scris în forma matriceală:
\[
  X' = AX, \quad A \in M_n(\mathbb{R}),
\]
unde $ A $ este matricea sistemului. Mai presupunem, de asemenea, că
matricea $ A $ este inversabilă, astfel încît sistemul admite soluția
$ X = 0 $.

Atunci avem:
\begin{theorem}[Poincar\'e-Liapunov]\label{thm:poi-li}
  Păstrînd contextul și notațiile de mai sus, dacă toate valorile proprii
  ale matricei $ A $ au partea reală negativă, atunci poziția de e\-chi\-li\-bru
  $ x = 0 $ este \emph{a\-simp\-to\-tic stabilă}.

  Dacă există $ \lambda \in \sigma(A) $, cu $\mathrm{Re}\lambda > 0 $, atunci
  $ x = 0 $ este instabilă.
\end{theorem}

%%%%%%%%%%%%%%%%%%%%%%%%%%%%%%%%%%%%%%%%%%%%%%%%%%%%%%%%%%%%%%%%%%%%%% 

\section{Exerciții}

1. Studiați stabilitatea poziției de echilibru $ x = 0 $ pentru sistemele
diferențiale:
\begin{enumerate}[(a)]
\item $ \begin{cases}
    x^\prime &= -4x + y \\
    y^\prime &= -x - 2y
  \end{cases} $;
\item $ \begin{cases}
    x^\prime &= -x + z \\
    y^\prime &= -2y - z \\
    z^\prime &= y - z
  \end{cases} $;
\item $ \begin{cases}
    x^\prime &= 2y \\
    y^\prime &= x + ay
  \end{cases}, a \in \RR $.
\end{enumerate}

\vspace{1cm}

2. Să se afle pentru ce valori ale lui $ a \in \RR $ soluția nulă a
sistemului de mai jos este asimptotic stabilă:
\[
  \begin{cases}
    x^\prime &= ax + y \\
    y^\prime &= (2 + a)x + ay \\
    z^\prime &= x + y - z
  \end{cases}
\]

%%%%%%%%%%%%%%%%%%%%%%%%%%%%%%%%%%%%%%%%%%%%%%%%%%%%%%%%%%%%%%%%%%%%%% 

\section{Integrale prime și linii de cîmp}

\begin{definition}\label{def:sist-autonom}
  \index{sistem!autonom}
  Un sistem diferențial de forma:
  \[
    x^\prime_j = v_j(x_1, \dots, x_n), \quad j = 1, \dots, n,
  \]
  unde $ v = (v_1, \dots, v_n) : U \to \RR^n $ este un cîmp de vectori definit
  într-un domeniu $ U \seq \RR^n $ se numește \emph{sistem diferențial autonom}.
\end{definition}

O altă formă de scriere a sistemului de mai sus este \emph{forma simetrică}:
\[
  \dfrac{dx_1}{v_1} = \dfrac{dx_2}{v_2} = \dots = \dfrac{dx_n}{v_n}.
\]
Dacă $ f : U \to \RR $ este o funcție oarecare, de clasă $ \kal{C}^1(U) $, atunci
pentru orice $ x \in U $ se poate scrie \emph{derivata lui $ f $ în $ x $ %
  în direcția vectorului $ v $}, notată $ \dfrac{df}{dv}(x) $, și definită
prin formula cunoscută:
\[
  \dfrac{df}{dv}(x) = \sum_{i = 1}^n \dfrac{\partial f}{\partial x_i} v_i(x).
\]

\begin{definition}\label{def:int-prim}
  \index{integrală!primă}
  \index{linii!de cîmp}
  Fie $ v : U \to \RR^n $ un cîmp de vectori și $ f : U \to \RR^n $ o funcție
  de clasă $ \kal{C}^1(U) $. Funcția se numește (o) \emph{integrală primă} a
  sistemului diferențial autonom $ x' = v(x), x \in U $ dacă derivata sa
  în direcția cîmpului de vectori $ v $ este nulă în fiecare punct din
  $ U $, adică $ \dfrac{df}{dv} = 0 $.
\end{definition}

Echivalent, putem înțelege definiția astfel: $ f : U \to \RR $ este o
integrală primă pentru sistemul diferențial autonom $ x' = v(x) $ dacă
și numai dacă pentru orice soluție $ x = \varphi(t), \varphi : I \to U $,
funcția $ f \circ \varphi $ este constantă pe $ I $. De aceea, uneori
mai putem scrie pe scurt $ f(x_1, \dots, x_n) = c $, constant. Rezultă
că putem înțelege integralele prime ca pe \emph{legi de conservare}.

În exerciții, pentru a găsi integralele prime asociate unui sistem autonom,
se rescrie sistemul în forma simetrică, iar apoi, folosind proprietățile rapoartelor
egale, încercăm să ajungem la un raport de forma $ \dfrac{df}{0} $, egal cu
rapoartele precedente. Atunci $ f $ va fi integrală primă, deoarece va
rezulta $ df = 0 $, adică $ f $ constantă în lungul curbelor integrale.

\textbf{Exemplu:} Să se găsească integralele prime ale sistemului simetric:
\[
  \dfrac{dx}{z - y} = \dfrac{dy}{x - z} = \dfrac{dz}{y - x}.
\]

\emph{Soluție:} Folosind proprietățile proporțiilor, rescriem sistemul:
\[
  \dfrac{dx}{z - y} = \dfrac{dy}{x - z} = \dfrac{dz}{y - x} = %
  \dfrac{dx + dy + dz}{z - x + x - z + y - x} = %
  \dfrac{dx + dy + dz}{0}.
\]

O altă formă de a obține o integrală primă este:
\[
  \dfrac{xdx + ydy + zdz}{x(z - x) + y(x - z) + z(y - x)} = %
  \dfrac{xdx + ydy + zdz}{0}.
\]

Rezultă:
\[
  \begin{cases}
    dx + dy + dz &= 0 \\
    xdx + ydy + zdz &= 0,
  \end{cases}
\]
de unde obținem două integrale prime:
\[
  x + y + z = c_1, \quad x^2 + y^2 + z^2 = c_2.
\]

\begin{remark}\label{rk:sist-cimp}
  Rezolvarea este identică în cazul în care cerința pornește cu un
  cîmp vectorial în spațiu, de forma:
  \[
    \vec{V} = (z - y)\vec{i} + (x - z)\vec{j} + (y - x) \vec{k}.
  \]
  Astfel, se asociază sistemul simetric de mai sus, etc.

  Pentru acest caz, integralele prime se mai numesc \emph{linii de cîmp}
  sau \emph{curbe integrale}.
\end{remark}

%%%%%%%%%%%%%%%%%%%%%%%%%%%%%%%%%%%%%%%%%%%%%%%%%%%%%%%%%%%%%%%%%%%%%% 

\section{Exerciții}

1. Să se determine liniile de cîmp pentru cîmpurile vectoriale de pe $\RR^3$:
\begin{enumerate}[(a)]
\item $\vec{v} = x\vec{i} + y \vec{j} + xyz \vec{k}$;
\item $\vec{v} = (2z - 3y)\vec{i} + (6x - 2z)\vec{j} + (3y - 6x)\vec{k}$;
\item $\vec{v} = (xz + y) \vec{i} + (x + yz) \vec{j} + (1 - z^2)\vec{k}$;
\item $\vec{v} = (x^2 + y^2)\vec{i} + 2xy \vec{j} - z^2 \vec{k}$.
\end{enumerate}

\emph{Indicație (a):} Scriem sistemul autonom asociat cîmpului vectorial $\vec{v}$
sub forma:
\[
  \frac{dx}{x} = \frac{dy}{y} = \frac{dz}{xyz}.
\]
Din prima egalitate, rezultă $x = c_1 y$, deci a doua egalitate devine:
\[
  c_1 ydy = \frac{dz}{z} \Rightarrow c_1 \frac{y^2}{2} = \ln |z| + c_2.
\]
Așadar, obținem $\dfrac{x}{y} \cdot \dfrac{y^2}{2} = \ln |z| + c_2$.

Rezultă că liniile de cîmp pentru $\vec{v}$ sînt curbele:
\[
  \begin{cases}
    \dfrac{x}{y} &= c_1 \\
    xy - \ln|z| &= c_2.
  \end{cases}
\]

\vspace{1cm}

(b) Sistemul poate fi scris sub forma:
\[
  \frac{dx}{2z - 3y} = \frac{dy}{6x - 2z} = \frac{dz}{3y-6x} = \frac{dx + dy + dz}{0}.
\]
Așadar, $dx + dy + dz = 0$, deci $x + y + z = c_1$.

Din forma inițială obținem și:
\[
  \frac{3xdx}{3x(2z - 3y)} = \frac{\frac{3}{2}ydy}{\frac{3}{2}y(3x - z)} = \frac{zdz}{z(y - 2x)} = %
  \frac{3xdx + \frac{3}{2}ydy + zdz}{0},
\]
deci $3xdx + \dfrac{3}{2}ydy + zdz = 0$, adică
$3\dfrac{x^2}{2} + \dfrac{3}{4}y^2 + \dfrac{z^2}{22} = c_2$.

\vspace{1cm}

(c) Obținem:
\[
  \frac{dx + dy}{(x + y)(z + 1)} = \frac{dz}{1 - z^2} \Rightarrow %
  \frac{d(x + y)}{x + y} = -\frac{dz}{z - 1}.
\]
Rezultă $\ln|x + y| + \ln|z - 1| = \ln c_1$, adică $(x + y)(z - 1) = c_1$.

Apoi:
\[
  \frac{dy - dx}{x + zy - xz - y} = \frac{dy - dx}{(x - y)(1 - z)} = \frac{dz}{1 - z^2}.
\]

Rezultă $-\dfrac{d(x - y)}{x - y} = \dfrac{dz}{z + 1}$, de unde $(x - y)(z + 1) = c_2$.

Așadar, liniile de cîmp sînt curbele:
\[
  \begin{cases}
    (x + y)(z - 1) &= c_1 \\
    (x - y)(z + 1) &= c_2
  \end{cases}
\]

\vspace{1cm}

(d) Sistemul simetric rezultat este:
\[
  \frac{dx}{x^2 + y^2} = \frac{dy}{2xy} = \frac{dz}{-z^2}.
\]

Prin adunarea primelor două rapoarte, obținem:
\[
  \frac{d(x + y)}{(x + y)^2} = \frac{dz}{-z^2},
\]
adică $\dfrac{1}{x + y} + \dfrac{1}{z} = c_1$.

Prin scădere, avem $\dfrac{d(x - y)}{(x - y)^2} = \dfrac{dz}{-z^2}$, deci
$\dfrac{1}{x - y} + \dfrac{1}{z} = c_2$.

Liniile de cîmp se obțin:
\[
  \begin{cases}
    \dfrac{1}{x + y} + \dfrac{1}{z} &= c_1 \\
    & \\
    \dfrac{1}{x - y} + \dfrac{1}{z} &= c_2
  \end{cases}
\]

%%%%%%%%%%%%%%%%%%%%%%%%%%%%%%%%%%%%%%%%%%%%%%%%%%%%%%%%%%%%%%%%%%%%%% 

\section{Resurse suplimentare}

\indent\indent Termenul în engleză pentru linii de cîmp este \texttt{field line}, iar cîteva
explicații, împreună cu exemple cunoscute din cazul cîmpurilor fizice, se pot
găsi pe \href{https://en.wikipedia.org/wiki/Field\_line}{Wikipedia}.

Metoda rezolvării ecuațiilor diferențiale în mod grafic, folosind \emph{cîmpul %
  tangent} (eng.\ \texttt{slope field}), cu semnificația fizică a cîmpului de
viteze, poate fi foarte utilă. Ea este prezentată succint pe
\href{https://en.wikipedia.org/wiki/Slope\_field}{Wikipedia} și mai detaliat,
inclusiv cu exerciții rezolvate, la \href{https://www.khanacademy.org/math/ap-calculus-ab/ab-differential-equations-new/ab-7-3/v/creating-a-slope-field}{Khan Academy}.

Trecerea de la cîmpuri vectoriale în spațiu la ecuații diferențiale (sau sisteme)
autonome este explicată pe scurt \href{http://www.math.tamu.edu/~zelenko/308\_w3\_qual.pdf}{aici}.


\vspace{1cm}

Alte materiale se mai pot găsi în cursul de la \href{http://www.mec.tuiasi.ro/diverse/matematici\_speciale.pdf}{Iași} (v.\ \S2.1.2),
în cursul \href{http://www.edumanager.ro/community/documente/ecuatii\_cu\_derivate\_partiale\_de\_ordinul\_al\_doilea.pdf}{prof.\ Crăciun și Barbu} (cap.\ 2)
și în \href{https://www.yumpu.com/ro/document/read/17874390/ecuatii-diferentiale-n-forma-simetrica-integrale-prime}{aceste notițe} de la UBB Cluj.

%%% Local Variables:
%%% mode: latex
%%% TeX-master: "../am2-18-19"
%%% End:
